\section{Chương~\ref{sec:nora}. \nt{NORA}}

Cấu trúc của hệ \nt{NORA}, hai chế độ của nó, mối liên hệ với đồng tử (không chỉ qua \nt{NORA}), sự tăng tỷ số tín hiệu trên nền và chữ U ngược trong hành vi đã được đo trực tiếp; mức khuếch đại nhân $g$ là mô hình; việc khuếch đại chuyển chế độ lựa chọn và hoạt động như nghịch đảo nhiệt độ là những điều suy ra trong chương; lợi ích của việc điều chỉnh khuếch đại là tính toán theo mô hình của sách; ``\nt{NORA} là núm vặn khám phá'' và ``\nt{NORA} là ngắt'' là giả thuyết.

{\footnotesize
\setlength{\tabcolsep}{3pt}\renewcommand{\arraystretch}{1.15}
\begin{longtable}{>{\raggedright}p{0.5cm}>{\raggedright}p{4.0cm}>{\raggedright}p{5.6cm}>{\raggedright}p{2.3cm}>{\raggedright\arraybackslash}p{1.7cm}}
\toprule
TT & Khẳng định & Thí nghiệm và điều đã đo & Nguồn & Trạng thái \\
\midrule\endfirsthead
\toprule
TT & Khẳng định & Thí nghiệm và điều đã đo & Nguồn & Trạng thái \\
\midrule\endhead
1 & Ở người có vài chục nghìn nơ-ron \nt{NORA}, gần như tất cả trong một nhóm & Lát cắt liên tiếp thân não người, nhuộm tyrosine hydroxylase: $53\,900$ tế bào trong nhân lục (locus coeruleus) và thêm $6\,260$ trong phân nhân lân cận. & \cite{Baker1989LC} & đã đo \\
2 & Đầu ra tỏa gần như khắp não, nhưng các phần của nhóm phần nào phục vụ những vùng khác nhau & Đánh dấu di truyền bằng virus ở chuột nhắt: nơ-ron \nt{NORA} của nhân lục nhận đầu vào từ những vùng rộng của não và gửi đầu ra tới những vùng rộng của não và tủy sống. Ở chuột cống: nhóm chiếu tới hạch hạnh nhân tăng cường việc học sợ hãi, còn một nhóm riêng chiếu tới vỏ não trước trán tăng cường việc dập tắt nó. & \cite{Schwarz2015LC, Uematsu2017LC, Poe2020} & đã đo \\
3 & Chế độ nền: xung đều với tần số từ dưới một tới vài xung mỗi giây, thay đổi theo mức tỉnh táo & Chuột cống, ghi nơ-ron nhân lục qua chu kỳ ngủ--thức: tần số cao nhất khi thức, thấp hơn trong giấc ngủ sóng chậm, gần bằng không trong giấc ngủ REM; thay đổi tần số đi trước sự chuyển trạng thái. & \cite{AstonJonesBloom1981} & đã đo \\
4 & Chế độ từng đợt: chùm xung ngắn với sự kiện quan trọng cho nhiệm vụ, vào khoảng thời điểm quyết định & Khỉ, nhiệm vụ tìm một mục tiêu hiếm: mọi nơ-ron nhân lục chỉ đáp ứng với mục tiêu bằng một chùm xung, $\sim90\ms$ sau mục tiêu và $\sim200\ms$ trước phản ứng bằng tay; khi làm kém, chùm xung yếu hơn. Trong nhiệm vụ chọn một trong hai, chùm xung gắn với phản ứng chặt hơn là với kích thích, và có cả khi trả lời đúng lẫn khi trả lời sai. & \cite{AstonJones1994vigilance, Clayton2004LC} & đã đo \\
5 & Đường kính đồng tử là điểm đo kiểm tra \nt{NORA} không chính xác: nó bám theo cả \nt{NORA}, cả \nt{ACET} và cả các vùng khác & Khỉ: xung của nhân lục, cho tới từng xung đơn lẻ của một tế bào, dự đoán sự giãn đồng tử; một mối liên hệ tương tự nhưng muộn hơn có ở các gò não và vỏ não đai. Chuột nhắt: dao động của đồng tử bám theo hoạt động của cả đầu ra noradrenaline lẫn đầu ra cholinergic trong vỏ não. & \cite{Joshi2016, Reimer2016} & đã đo \\
6 & Noradrenaline ức chế hoạt động nền mạnh hơn hoạt động được gợi ra; mức khuếch đại nhân~\eqref{eq:gain} là mô hình thể hiện hiệu ứng này & Khỉ tỉnh, vỏ não thính giác, điện di ion noradrenaline: hoạt động nền của nơ-ron bị ức chế mạnh hơn đáp ứng với tiếng kêu của đồng loại --- tỷ số tín hiệu trên nhiễu tăng. Hàm sigmoid nhân~\eqref{eq:gain} lấy từ một mô hình mạng; việc độ dốc bị nhân lên theo đúng nghĩa đen chưa được đo. & \cite{Foote1975NE, ServanSchreiber1990} & đã đo; $g$ là mô hình \\
7 & Khuếch đại chỉ làm tăng $d'$ trước nhiễu sau bộ khuếch đại~\eqref{eq:gainsnr} (mệnh đề~\ref{prop:gainnoise}) & Suy ra trong chương; trong mô hình mạng, khuếch đại cải thiện việc phát hiện tín hiệu. Chưa có thí nghiệm tách nhiễu trước và sau bộ khuếch đại và kiểm tra sự khác biệt này. & suy ra trong chương; \cite{ServanSchreiber1990} & lý thuyết \\
8 & Ranh giới $g\beta=1$ chuyển mạng lựa chọn từ ``cân nhắc'' sang ``đã quyết''~\eqref{eq:wtagain} (mệnh đề~\ref{prop:gainwta}, hình~\ref{fig:pitchfork}) & Phân tích tuyến tính hai nhóm ức chế lẫn nhau; suy ra trong chương. Chưa có phép đo trực tiếp ngưỡng này trong não. & suy ra trong chương & lý thuyết \\
9 & Hệ số khuếch đại là nghịch đảo nhiệt độ của lựa chọn~\eqref{eq:softmax} & Với nhiễu logistic sau bộ khuếch đại, xác suất lựa chọn là softmax với $T=\sigma/g$; suy ra trong chương. & suy ra trong chương & lý thuyết \\
10 & \nt{NORA} đặt mức khám phá: hoạt động nền cao --- $g$ hiệu dụng nhỏ (khám phá), chùm xung từng đợt lúc quyết định --- lớn (quyết định) & Người: trước một lựa chọn ngẫu nhiên, đồng tử nền rộng hơn so với trước khi chọn phương án tốt nhất đã biết; người có đồng tử rộng hơn thiên về khám phá hơn. Chuột cống: việc chuyển sang chế độ chọn ngẫu nhiên được điều khiển bởi đầu vào \nt{NORA} tới vỏ não đai trước. Việc chùm xung làm tăng $g$ hiệu dụng chưa được đo trực tiếp. & \cite{JepmaNieuwenhuis2011, Tervo2014stochastic, AstonJones2005} & giả thuyết \\
11 & Phần thưởng phụ thuộc vào $g$ theo chữ U ngược; $g$ tốt nhất nhỏ hơn trong thế giới thay đổi nhanh (mệnh đề~\ref{prop:explore}, hình~\ref{fig:invertedU}) & \texttt{models/nora\_explore.py}, $60$ lần chạy, mỗi lần $4000$ lượt thử. Thay đổi một lần sau mỗi $1000$ lượt: $g=16$ tốt nhất, tỷ lệ $0.976$; sau mỗi $20$ lượt: $g=8$, tỷ lệ $0.886$; khi $g=0.5$ --- $0.77$. Tính lại khớp với chương. & \texttt{models/\allowbreak nora\_\allowbreak explore.py} & mô hình \\
12 & Chữ U ngược trong hành vi: tốt nhất khi hoạt động nền vừa phải với chùm xung rõ nét & Khỉ, cùng nhiệm vụ như ở dòng~4: trong những giai đoạn làm tốt, tần số nền vừa phải và chùm xung với mục tiêu rõ nét; khi tần số nền cao thì không có chùm xung và nhiều phản ứng sai hơn. Thay đổi của hoạt động nền và hoạt động được gợi ra bám theo dao động của chất lượng thực hiện. & \cite{Usher1999LC, AstonJones2005} & đã đo \\
13 & \nt{NORA} báo sự bất định không được dự liệu, \nt{ACET} --- sự bất định được dự liệu & Lý thuyết suy luận Bayes với hai loại bất định; các công trình được trích không có thí nghiệm trực tiếp nào tách các tín hiệu này theo chất dẫn truyền. & \cite{YuDayan2005} & giả thuyết \\
14 & Sau một thay đổi đột ngột, người ta học nhanh hơn, và đồng tử giãn ra & Người, nhiệm vụ dự đoán với những thay đổi bất ngờ: những thay đổi ngắn của đồng tử bám theo ước lượng xác suất có thay đổi, và cùng với đồng tử nền dự đoán dữ liệu mới làm thay đổi ước lượng mạnh tới mức nào (tốc độ học); thay đổi đồng tử không liên quan tới ánh sáng và nhiệm vụ làm thay đổi tốc độ này. Việc ở đây đồng tử phản ánh chính \nt{NORA} là suy luận từ dữ liệu gián tiếp ở dòng~5. & \cite{Nassar2012} & đã đo \\
15 & Chùm xung từng đợt hoạt động như một ngắt: mạng xóa trạng thái & Chưa có thí nghiệm trực tiếp nào trong đó chùm xung \nt{NORA} xóa trạng thái của mạng; mới chỉ đo được chùm xung với các sự kiện quan trọng (dòng~4). & \cite{BouretSara2005} & giả thuyết \\
16 & Điều chỉnh $g$ hoặc tốc độ học theo sự bất ngờ hầu như không tốt hơn các thiết lập cố định tốt nhất & \texttt{models/nora\_explore.py}, thế giới với các thời kỳ dài $1000$ lượt thử (thay đổi một lần sau mỗi $1000$ và sau mỗi $20$ lượt). tốt nhất với $g$ cố định là $g=8$: $0.925$; điều chỉnh $g$: tới $0.927$; điều chỉnh tốc độ học: tới $0.926$; tốc độ cố định $0.4$: $0.928$. Tính lại khớp với chương. & \texttt{models/\allowbreak nora\_\allowbreak explore.py} & mô hình \\
\bottomrule
\end{longtable}}
